\section{Geometric Tightness and Continuum Limit}
\label{app:tightness_continuum}
\label{sec:continuum-details}

\textbf{Note: This section has been revised in accordance with the Rigorization Addendum (Appendix \ref{app:rigor-addendum}).}

This appendix formulates the precise condition required for the existence of the continuum limit. Instead of relying on heuristic scaling arguments or unproven plaquette moment bounds, we state the tightness requirement as a rigorous hypothesis on the distribution of gauge fields.

\subsection{Hypothesis: Tightness in Negative Sobolev Space}

The existence of a non-trivial continuum limit is guaranteed if the sequence of lattice measures satisfies the following tightness condition.

\begin{hypothesis}[Tightness in $H^{-s}$]
\label{hyp:tightness-sobolev-v2}
Let $\mu_a$ be the probability measure on the space of lattice gauge fields at spacing $a$. We embed these fields into the space of distributions on $\mathbb{R}^4$ (or $\mathbb{T}^4$) via a suitable interpolation map $\iota_a: \mathcal{A}_{lat} \to \mathcal{D}'(\mathbb{R}^4)$.

We assume there exists an $s > 2$ and $p \ge 1$ such that:
\begin{equation}
\sup_{a \to 0} \mathbb{E}_{\mu_a} \left[ \| \iota_a(A) \|_{H^{-s}}^p \right] < \infty
\end{equation}
where $H^{-s}$ is the negative Sobolev space of index $s$.

Furthermore, we assume a "log-chart" large-field suppression:
\begin{equation}
\mathbb{P}_{\mu_a} \left( \sup_{x} |A(x)| > \Lambda \right) \le C e^{-c \Lambda^2}
\end{equation}
for the smoothed field at scale $a$.
\end{hypothesis}

\subsection{Consequences of the Hypothesis}

If Hypothesis \ref{hyp:tightness-sobolev-v2} holds, then by Prokhorov's Theorem, the family of measures $\{\mu_a\}$ is pre-compact in the weak topology of measures on $H^{-s}$.

\begin{theorem}[Existence of Continuum Limit]
Assume Hypothesis \ref{hyp:tightness-sobolev-v2}. Then there exists a subsequence $a_k \to 0$ and a probability measure $\mu$ on $H^{-s}(\mathbb{R}^4)$ such that $\mu_{a_k} \to \mu$ weakly.
\end{theorem}

This limit measure $\mu$ is the candidate for the continuum Yang-Mills measure. The subsequent task (Step 3 in the Rigorous Proof Chain) is to verify that $\mu$ satisfies the Osterwalder-Schrader axioms, which follows from the RG-Stability hypothesis.

\subsection{Why this replaces the $a^{2p}$ assumption}

The previous draft assumed $\mathbb{E}[|\mathcal{P}(U_p)|^p] \le C_p a^{2p}$. While heuristically plausible from asymptotic freedom, this bound is difficult to prove directly without full control of the renormalization group. The Sobolev tightness hypothesis is a weaker and more natural condition for distribution-valued fields, and it is the standard setting for Constructive QFT (e.g., in the $P(\phi)_2$ construction).

\begin{remark}
Proving Hypothesis \ref{hyp:tightness-sobolev-v2} requires establishing that the renormalized effective actions do not develop singularities that grow too fast in the UV limit. This is the content of the "UV Regularity" problem.
Full rigorous derivation of the conditions for this hypothesis from RG stability is provided in Appendix \ref{app:rigorous_rg_tightness}.
\end{remark}

